Las instituciones financieras y los organismos encargados de la política económica deberían mantener un monitoreo permanente de indicadores como las remesas familiares, el numerario en circulación y los medios de pago, debido a su estrecha relación con la evolución de la actividad económica nacional.
Para estudios orientados a la formulación de políticas públicas o decisiones estratégicas, se recomienda complementar este análisis descriptivo con modelos econométricos que permitan cuantificar el impacto y la causalidad entre las variables analizadas.
Es recomendable utilizar información estadística actualizada y revisar periódicamente las series publicadas por el Banco de Guatemala, ya que las cifras más recientes pueden ser objeto de revisiones o ajustes metodológicos.
Finalmente, resulta conveniente impulsar políticas que fortalezcan la inclusión financiera y promuevan una mayor bancarización, con el propósito de reducir la dependencia del efectivo y canalizar una mayor proporción de los recursos provenientes de remesas hacia el sistema financiero formal.